\section{User define quantity MOC conventions\label{sec:udmoc1}}

The conventions adopted for FMOCs are generic and could be applied to any quantity
$q$ numerically represented as a 64-bit floating-point value.
A user wishing to define a custom QMOC (where $Q$ stands for Quantity) would only need to specify:
\begin{itemize}
  \item the validity range of the quantity: $[q_{min}, q_{max}[$;
  \item the mapping function: $g_q: [q_{min}, q_{max}[ \rightarrow [0,1[$.
\end{itemize}

However, one of the advantages of MOCs is the use of standardized quantities expressed
in well-defined units and reference systems.
In order to perform efficient and meaningful set operations between two MOCs, both MOCs must
be defined using the same quantity, expressed in the same units, and using
the same validity interval and mapping function.

Poorly constrained or insufficiently documented user-defined quantities may therefore compromise the interoperability of MOCs.
